% =====================================================================
%  Chương 2 — Hệ thống Chẩn đoán Nguy cơ Tiểu đường
% =====================================================================
\chapter{Hệ thống Chẩn đoán Nguy cơ Tiểu đường}
\label{ch:tieuduong}

\section{Phát biểu bài toán}
\label{sec:phatbieu-dia}

\subsection{Hiện tượng thực tế và người dùng}

Đái tháo đường type 2 tiến triển âm thầm: người bệnh thường không có triệu
chứng rõ rệt cho tới khi biến chứng đã hình thành. Vì vậy giá trị của một hệ
thống trong miền này không nằm ở chỗ ``chẩn đoán chính xác'' --- việc ấy thuộc
về xét nghiệm HbA1c tĩnh mạch --- mà nằm ở chỗ \textbf{sàng lọc}: từ những chỉ
số đo được rẻ và nhanh, chỉ ra ai cần đi làm xét nghiệm khẳng định.

Việc xác định đúng vai trò này quyết định mọi lựa chọn kỹ thuật về sau. Nó là
lý do báo cáo chấm điểm bằng $F_1$ và Recall chứ không bằng Accuracy
(Mục~\ref{sec:baseline}), và là lý do Thực nghiệm~2 hạ ngưỡng quyết định xuống
tận $0{,}15$ (Mục~\ref{sec:tn2-dia}).

\subsection{Phát biểu hình thức}

Cho một hồ sơ y sinh $\mathbf{x}\in\mathbb{R}^{d}$ gồm các chỉ số đo được của
một người, hãy học hàm

\begin{equation}
f_{\bm\theta}:\ \mathbb{R}^{d}\longrightarrow[0,1],
\qquad
\hat p = f_{\bm\theta}(\mathbf{x}) = P\big(\texttt{Outcome}=1 \mid \mathbf{x}\big),
\label{eq:baitoan-dia}
\end{equation}

\noindent rồi quy xác suất về nhãn nhị phân qua một ngưỡng quyết định $\theta$:

\begin{equation}
\hat y = \mathbb{1}\big[\,\hat p \ge \theta\,\big],
\qquad \theta \in (0,1).
\label{eq:nguong-dia}
\end{equation}

Ngưỡng $\theta$ ở~\eqref{eq:nguong-dia} là một tham số \textbf{của hệ thống},
không phải của mô hình: nó không được học từ dữ liệu mà được chọn theo mục tiêu
sử dụng. Việc tách bạch hai thứ này là điều kiện cần để Thực nghiệm~2 có nghĩa.

\begin{table}[H]
\centering
\caption{Đặc tả bài toán chẩn đoán nguy cơ tiểu đường}
\label{tab:dacta-dia}
\small
\begin{tabularx}{\textwidth}{@{}L{4.2cm}X@{}}
\toprule
\textbf{Hạng mục} & \textbf{Nội dung} \\
\midrule
Hiện tượng thực tế   & Tình trạng chuyển hoá glucose và nguy cơ mắc đái tháo đường type 2 \\
Một quan sát         & Một hồ sơ y sinh của một người trưởng thành \\
Loại bài toán        & Phân lớp nhị phân (Binary Classification) \\
Biến mục tiêu        & \texttt{Outcome} $\in\{0,1\}$ \\
Số quan sát          & 768 hồ sơ \\
Đặc trưng đầu vào    & 8 cột gốc, tất cả đều là số; không có biến định danh \\
Đặc trưng sau thiết kế & 12 cột (8 gốc $+$ 4 phái sinh) \\
Người dùng cuối      & Bác sĩ sàng lọc, và chính người được sàng lọc \\
Độ đo quyết định     & $F_1$ và Recall --- không phải Accuracy \\
\bottomrule
\end{tabularx}
\end{table}

\section{Tập dữ liệu và nguồn gốc}
\label{sec:dulieu-dia}

Tập \textbf{Pima Indians Diabetes}~\cite{pima} do Viện Quốc gia về Đái tháo đường, Bệnh
Tiêu hoá và Bệnh Thận Hoa Kỳ thu thập, công bố công khai trên Kaggle:

\begin{center}
\small\url{https://www.kaggle.com/datasets/uciml/pima-indians-diabetes-database}
\end{center}

Toàn bộ đối tượng là nữ, từ 21 tuổi trở lên, thuộc cộng đồng người Pima --- một
giới hạn về phạm vi áp dụng sẽ được nêu lại ở Mục~\ref{sec:hanche}.

\begin{table}[H]
\centering
\caption{Đặc tả tám đặc trưng gốc của tập Pima Indians Diabetes}
\label{tab:dactrung-dia}
\small
\begin{tabularx}{\textwidth}{@{}L{3.5cm}L{1.6cm}R{1.5cm}R{1.5cm}X@{}}
\toprule
\textbf{Đặc trưng} & \textbf{Đơn vị} & \textbf{Trung bình} & \textbf{Trung vị} & \textbf{Ý nghĩa} \\
\midrule
Pregnancies              & lần   & 3{,}845   & 3{,}0   & Số lần mang thai \\
Glucose                  & mg/dL & 121{,}656 & 117{,}0 & Đường huyết sau nghiệm pháp dung nạp 2 giờ \\
BloodPressure            & mmHg  & 72{,}387  & 72{,}0  & Huyết áp tâm trương \\
SkinThickness            & mm    & 29{,}108  & 29{,}0  & Độ dày nếp gấp da cơ tam đầu \\
Insulin                  & $\mu$U/mL & 140{,}672 & 125{,}0 & Insulin huyết thanh sau 2 giờ \\
BMI                      & kg/m$^2$ & 32{,}455 & 32{,}3 & Chỉ số khối cơ thể \\
DiabetesPedigreeFunction & --    & 0{,}472   & 0{,}372 & Chỉ số di truyền gia đình \\
Age                      & năm   & 33{,}241  & 29{,}0  & Tuổi \\
\bottomrule
\end{tabularx}
\end{table}

\nhanxet{Trung bình lớn hơn trung vị ở hầu hết các cột --- dấu hiệu của phân
phối lệch phải. Riêng \texttt{Age} chênh nhiều nhất ($33{,}2$ so với $29{,}0$),
cho thấy đuôi phải kéo dài: mẫu chủ yếu là người trẻ, kèm một nhóm nhỏ tuổi
cao. Các con số trong bảng được tính \textbf{sau} khi áp quy ước ``0 là khuyết''
ở Mục~\ref{sec:lamsach-dia}; tính trên dữ liệu thô sẽ cho trung bình
\texttt{Insulin} thấp hơn thực tế gần một nửa.}

\section{Làm sạch dữ liệu: khi số 0 không phải là số đo}
\label{sec:lamsach-dia}

Tập dữ liệu không có ô trống nào theo nghĩa kỹ thuật --- mọi ô đều có giá trị.
Nhưng năm cột sinh học chứa giá trị $0$, và trên một người sống thì nồng độ
glucose bằng $0$, huyết áp bằng $0$ hay BMI bằng $0$ đều là điều bất khả. Đây
là \textbf{quy ước mã hoá cho ``không đo được''}, không phải số đo thật.

\begin{table}[H]
\centering
\caption{Số quan sát mang giá trị 0 phi lý ở năm cột sinh học}
\label{tab:zero-dia}
\small
\begin{tabular}{@{}lrrl@{}}
\toprule
\textbf{Cột} & \textbf{Số ca} & \textbf{Tỷ lệ} & \textbf{Vì sao 0 là bất khả} \\
\midrule
Insulin       & 374 & 48{,}7\% & Cơ thể luôn tiết insulin nền \\
SkinThickness & 227 & 29{,}6\% & Da luôn có độ dày đo được \\
BloodPressure &  35 &  4{,}6\% & Huyết áp 0 nghĩa là ngừng tuần hoàn \\
BMI           &  11 &  1{,}4\% & BMI 0 nghĩa là khối lượng cơ thể bằng 0 \\
Glucose       &   5 &  0{,}7\% & Glucose 0 không tương thích với sự sống \\
\bottomrule
\end{tabular}
\end{table}

\begin{figure}[H]
\centering
\includegraphics[width=0.78\textwidth]{dia_missing_ratio.png}
\caption[Tỷ lệ dữ liệu khuyết sau khi áp quy ước ``0 là khuyết'']%
{Tỷ lệ dữ liệu khuyết ở năm cột sinh học sau khi áp quy ước ``giá trị 0 là dữ
liệu không đo được''}
\label{fig:dia-missing}
\end{figure}

\bieudo{Biểu đồ cột ngang (horizontal bar chart), sắp xếp giảm dần theo tỷ lệ
khuyết. Trục ngang là tỷ lệ phần trăm, trục dọc là tên cột.}

\phanphoi{Phân bố cực kỳ không đều giữa các cột. Hai cột đầu chiếm gần như toàn
bộ lượng khuyết, năm cột còn lại gộp lại chưa tới $7\%$.}

\chitiet{\texttt{Insulin} khuyết $48{,}7\%$ --- gần một nửa dữ liệu.
\texttt{SkinThickness} khuyết $29{,}6\%$. Ba cột còn lại lần lượt $4{,}6\%$,
$1{,}4\%$ và $0{,}7\%$. Không cột nào vượt ngưỡng loại bỏ $80\%$ nên cả năm cột
đều được giữ lại và điền khuyết bằng trung vị.}

\begin{phathien}{Rò rỉ nhãn qua bước điền khuyết --- một sai lầm thật đã xảy ra}
\label{sec:roro}
Bản triển khai đầu tiên điền khuyết bằng \emph{trung vị theo nhóm} \texttt{Outcome}
--- tức người bệnh được điền bằng trung vị của người bệnh. Kết quả nhảy lên
Accuracy $0{,}8705$, vượt xa trần $0{,}77$--$0{,}79$ mà tài liệu công bố trên
cùng tập dữ liệu này~\cite{smith1988}. Con số đẹp bất thường ấy chính là dấu hiệu của lỗi.

Nguyên nhân: với \texttt{Insulin} khuyết $374/768$ và \texttt{SkinThickness}
khuyết $227/768$, gần một nửa dữ liệu sau khi điền đã mang dấu vết của chính
nhãn cần dự đoán. Mô hình không học quan hệ sinh học; nó đọc ngược ra nhãn từ
giá trị đã điền.

Cách sửa: đưa \texttt{SimpleImputer(strategy="median")} vào \emph{bên trong}
\texttt{Pipeline}, để \texttt{fit} chỉ nhìn thấy tập huấn luyện. Accuracy trở
về $0{,}7914$ --- thấp hơn, nhưng là con số thật. \textbf{Bài học: một kết quả
vượt trần tài liệu là lý do để đi tìm lỗi, không phải lý do để ăn mừng.}
\end{phathien}

\section{Phân tích khám phá dữ liệu}
\label{sec:eda-dia}

\subsection{Phân phối biến mục tiêu}

\begin{figure}[H]
\centering
\includegraphics[width=0.56\textwidth]{dia_pie_outcome.png}
\caption[Phân phối hai lớp của biến mục tiêu \texttt{Outcome}]%
{Phân phối hai lớp của biến mục tiêu \texttt{Outcome}}
\label{fig:dia-pie}
\end{figure}

\bieudo{Biểu đồ tròn (pie chart) --- dạng phù hợp khi cần thể hiện tỷ trọng của
một số ít hạng mục trên tổng thể.}

\phanphoi{Lệch lớp ở mức trung bình, tỷ lệ xấp xỉ $2:1$ nghiêng về lớp âm tính.}

\chitiet{500 ca âm tính chiếm $65{,}1\%$ và 268 ca dương tính chiếm $34{,}9\%$.
Mức lệch này chưa tới ngưỡng ``mất cân bằng nghiêm trọng'' (thường tính từ
$10:1$) nên không cần lấy mẫu lại, nhưng đã đủ để làm Accuracy mất ý nghĩa:
một mô hình đoán bừa ``tất cả đều khoẻ'' đạt ngay $65{,}1\%$.}

\subsection{Phân phối đặc trưng quan trọng nhất}

\begin{figure}[H]
\centering
\includegraphics[width=0.78\textwidth]{dia_hist_glucose.png}
\caption[Phân phối tần suất của đặc trưng \texttt{Glucose}]%
{Phân phối tần suất của đặc trưng \texttt{Glucose}}
\label{fig:dia-hist-glucose}
\end{figure}

\bieudo{Biểu đồ tần suất (histogram). Dải giá trị liên tục được chia thành các
khoảng bằng nhau, chiều cao mỗi cột là số quan sát rơi vào khoảng đó.}

\phanphoi{Dạng gần chuẩn nhưng lệch phải nhẹ, một đỉnh duy nhất (unimodal).}

\chitiet{Khối lượng dữ liệu tập trung trong khoảng $90$--$140$ mg/dL, đỉnh tại
khoảng $105$--$115$ mg/dL với gần 70 quan sát. Đuôi trái mỏng, bắt đầu từ
khoảng $45$ mg/dL. Đuôi phải dày hơn hẳn và kéo tới $199$ mg/dL --- vùng này
tương ứng với ngưỡng chẩn đoán đái tháo đường ($\ge 200$ mg/dL sau nghiệm pháp
dung nạp), nên độ dày của đuôi phải chính là tín hiệu mà mô hình khai thác.}

\subsection{So sánh phân phối giữa hai nhóm}

\begin{figure}[H]
\centering
\includegraphics[width=\textwidth]{dia_kde_by_class.png}
\caption[Phân phối mật độ ba đặc trưng chính, tách theo tình trạng bệnh]%
{Ước lượng mật độ hạt nhân (KDE) của \texttt{Glucose}, \texttt{BMI} và
\texttt{Age}, tách theo tình trạng bệnh}
\label{fig:dia-kde-class}
\end{figure}

\bieudo{Ba biểu đồ mật độ hạt nhân chồng lớp (overlaid KDE). Khác biểu đồ tần
suất, KDE dùng đường cong liên tục nên hình dáng phân phối không phụ thuộc vào
cách chọn độ rộng khoảng. Hai màu là hai nhóm của biến mục tiêu.}

\phanphoi{Cả ba cặp đường cong đều lệch nhau theo cùng một hướng --- nhóm mắc
bệnh dịch sang phải --- nhưng \emph{mức độ} tách rời rất khác nhau.}

\chitiet{\texttt{Glucose} tách rõ nhất: đỉnh nhóm khoẻ ở khoảng $105$ mg/dL,
đỉnh nhóm bệnh ở khoảng $135$--$145$ mg/dL, và phần chồng lấn tuy còn lớn nhưng
hai đỉnh đã tách hẳn. \texttt{BMI} tách vừa phải, hai đỉnh chỉ cách nhau vài đơn
vị quanh mốc $30$--$35$ kg/m$^2$. \texttt{Age} tách yếu nhất: nhóm khoẻ có một
đỉnh nhọn ở khoảng $22$--$25$ tuổi, nhóm bệnh phân bố thoải và rộng hơn nhưng
phần chồng lấn chiếm gần trọn vùng dữ liệu.}

\nhanxet{Thứ tự tách rời ở đây --- \texttt{Glucose} $>$ \texttt{BMI} $>$
\texttt{Age} --- báo trước đúng thứ tự tương quan với nhãn ở
Bảng~\ref{tab:corr-dia} và đúng thứ tự mức quan trọng đặc trưng ở
Hình~\ref{fig:dia-importance}. Ba phương pháp độc lập cùng chỉ về một kết luận
là bằng chứng mạnh hơn bất kỳ phương pháp nào đứng riêng.}

\subsection{Giá trị ngoại lai}

\begin{figure}[H]
\centering
\includegraphics[width=\textwidth]{dia_box_grid.png}
\caption[Biểu đồ hộp của tám đặc trưng định lượng]%
{Biểu đồ hộp của tám đặc trưng định lượng}
\label{fig:dia-box}
\end{figure}

\bieudo{Lưới tám biểu đồ hộp (boxplot). Hộp bao phủ khoảng tứ phân vị Q1--Q3,
đường giữa là trung vị, râu kéo dài tối đa $1{,}5\times$IQR, mọi điểm ngoài râu
được đánh dấu là ngoại lai.}

\phanphoi{Ba nhóm khác hẳn nhau. Nhóm gọn gàng: \texttt{Glucose},
\texttt{BloodPressure}, \texttt{BMI} --- hộp cân đối, ít điểm ngoài râu. Nhóm
lệch phải mạnh: \texttt{Insulin}, \texttt{DiabetesPedigreeFunction},
\texttt{Age}. Nhóm đặc biệt: \texttt{SkinThickness}, có hộp rất hẹp.}

\chitiet{\texttt{Insulin} có hộp nén sát đáy với một dải điểm ngoại lai dày đặc
kéo lên tới hơn $800$ $\mu$U/mL --- hệ quả trực tiếp của việc $48{,}7\%$ giá trị
được điền bằng cùng một con số trung vị, khiến hộp co lại quanh giá trị ấy.
\texttt{SkinThickness} cũng cho hình dáng tương tự vì lý do y hệt.
\texttt{Age} có hộp trải từ khoảng $25$ tới $41$ tuổi với đuôi trên kéo tới
$81$. \texttt{Pregnancies} là biến đếm nên hình dáng bậc thang, tối đa $17$ lần.}

\nhanxet{Không giá trị ngoại lai nào bị loại bỏ. Với dữ liệu y sinh, một chỉ số
cực đoan thường là \emph{ca bệnh nặng} --- đúng nhóm mà hệ thống sàng lọc cần
phát hiện nhất. Loại bỏ chúng là loại bỏ tín hiệu, không phải loại bỏ nhiễu.}

\section{Thiết kế đặc trưng}
\label{sec:dactrung-dia}

Bốn đặc trưng phái sinh được xây từ tri thức y học, không phải từ tìm kiếm mù:

\begin{table}[H]
\centering
\caption{Bốn đặc trưng thiết kế và căn cứ y học}
\label{tab:engineered-dia}
\small
\begin{tabularx}{\textwidth}{@{}L{3.2cm}L{4.0cm}X@{}}
\toprule
\textbf{Đặc trưng} & \textbf{Công thức} & \textbf{Căn cứ} \\
\midrule
\texttt{Glucose\_BMI\_Risk} & $\texttt{Glucose}\times\texttt{BMI}$ &
Tăng đường huyết đi kèm béo phì có rủi ro cao hơn tổng rủi ro của hai yếu tố
riêng lẻ --- một tương tác mà mô hình tuyến tính không tự tạo ra được \\
\addlinespace
\texttt{Glucose\_Level} & Rời rạc hoá theo ngưỡng lâm sàng &
Ba mức: bình thường, tiền đái tháo đường, đái tháo đường \\
\addlinespace
\texttt{BMI\_Class} & Rời rạc hoá theo phân loại WHO &
Bốn mức: gầy, bình thường, thừa cân, béo phì \\
\addlinespace
\texttt{Age\_Group} & Rời rạc hoá theo nhóm tuổi &
Rủi ro tăng không tuyến tính theo tuổi \\
\bottomrule
\end{tabularx}
\end{table}

\section{Ma trận tương quan và sàng lọc đặc trưng}
\label{sec:corr-dia}

\begin{figure}[H]
\centering
\includegraphics[width=\textwidth]{dia_corr_heatmap.png}
\caption[Ma trận tương quan Pearson của 8 đặc trưng gốc và 4 đặc trưng thiết kế]%
{Ma trận tương quan Pearson --- 8 đặc trưng gốc, 4 đặc trưng thiết kế và biến
mục tiêu}
\label{fig:dia-corr}
\end{figure}

\bieudo{Bản đồ nhiệt (heatmap) đối xứng. Màu đỏ là tương quan dương, màu xanh là
tương quan âm, độ đậm tỷ lệ với độ lớn hệ số.}

\phanphoi{Toàn bộ ma trận nghiêng về phía đỏ nhạt --- không có cặp nào tương
quan âm đáng kể. Ba cụm đậm nổi lên rõ, tất cả đều là cặp gốc -- phái sinh.}

\chitiet{Ba cặp đậm nhất đều do định nghĩa mà có:
\texttt{Glucose}--\texttt{Glucose\_Level} $=0{,}87$,
\texttt{Glucose}--\texttt{Glucose\_BMI\_Risk} $=0{,}82$,
\texttt{Age}--\texttt{Age\_Group} $=0{,}95$ và
\texttt{BMI}--\texttt{BMI\_Class} $=0{,}78$. Với biến mục tiêu, cột
\texttt{Outcome} có màu nhạt hơn hẳn: hệ số cao nhất chỉ là $0{,}52$.}

\begin{table}[H]
\centering
\caption{Tương quan Pearson với biến mục tiêu \texttt{Outcome}, sắp giảm dần}
\label{tab:corr-dia}
\small
\begin{tabular}{@{}lrlr@{}}
\toprule
\textbf{Đặc trưng} & \textbf{$\rho$} & \textbf{Đặc trưng} & \textbf{$\rho$} \\
\midrule
\texttt{Glucose\_BMI\_Risk} & \textbf{0{,}5199} & \texttt{Age\_Group}     & 0{,}2322 \\
\texttt{Glucose}            & 0{,}4928 & \texttt{Pregnancies}    & 0{,}2219 \\
\texttt{Glucose\_Level}     & 0{,}4356 & \texttt{SkinThickness}  & 0{,}2149 \\
\texttt{BMI}                & 0{,}3120 & \texttt{Insulin}        & 0{,}2038 \\
\texttt{BMI\_Class}         & 0{,}3098 & \texttt{DiabetesPedigreeFunction} & 0{,}1738 \\
\texttt{Age}                & 0{,}2384 & \texttt{BloodPressure}  & 0{,}1657 \\
\bottomrule
\end{tabular}
\end{table}

\nhanxet{Đặc trưng thiết kế \texttt{Glucose\_BMI\_Risk} tương quan với nhãn
\emph{mạnh hơn} cả \texttt{Glucose} gốc ($0{,}5199$ so với $0{,}4928$) ---
bằng chứng định lượng cho thấy tri thức y học đưa vào đúng chỗ có giá trị. Ngưỡng
sàng lọc $|\rho|\ge0{,}20$ giữ lại năm đặc trưng gốc: \texttt{Glucose},
\texttt{SkinThickness}, \texttt{Insulin}, \texttt{BMI}, \texttt{Age}.

Cần nói rõ giới hạn: hệ số Pearson chỉ đo quan hệ \textbf{tuyến tính}. Giá trị
gần $0$ không loại trừ quan hệ phi tuyến mạnh, nên bảng này được dùng để
\emph{xếp hạng} chứ không dùng làm căn cứ duy nhất để loại đặc trưng --- mọi mô
hình trong chương vẫn nhận đủ 12 cột.}

\section{Chia dữ liệu và mô hình cơ sở}

Áp dụng giao thức ba tập ở Mục~\ref{sec:giaothuc}, có phân tầng theo nhãn để
giữ nguyên tỷ lệ $65/35$ trong cả ba tập:

\begin{table}[H]
\centering
\caption{Kích thước ba tập dữ liệu và kết quả mô hình cơ sở}
\label{tab:split-dia}
\small
\begin{tabular}{@{}lrr@{\hspace{1.5cm}}lr@{}}
\toprule
\textbf{Tập} & \textbf{Số ca} & \textbf{Tỷ lệ} & \textbf{Mô hình cơ sở} & \textbf{Giá trị} \\
\midrule
Huấn luyện & 552 & 71{,}9\% & Accuracy  & 0{,}6547 \\
Kiểm tra   & 139 & 18{,}1\% & Precision & 0{,}0000 \\
Ngoại kiểm &  77 & 10{,}0\% & Recall    & \textbf{0{,}0000} \\
\textbf{Tổng} & \textbf{768} & \textbf{100\%} & $F_1$ & 0{,}0000 \\
\bottomrule
\end{tabular}
\end{table}

\begin{phathien}{Vì sao Accuracy $0{,}6547$ là một con số nguy hiểm}
Mô hình cơ sở luôn đoán ``khoẻ mạnh''. Nó đúng $65{,}47\%$ số lần --- nghe không
tệ. Nhưng Recall bằng $0$: nó \textbf{bỏ sót toàn bộ 48 người bệnh} trong tập
kiểm tra. Một hệ thống sàng lọc như vậy vô dụng hoàn toàn, dù chỉ số Accuracy
của nó cao hơn cả XGBoost ($0{,}7194$)~\dots{} không, thực ra thấp hơn --- nhưng
khoảng cách chỉ là $6$ điểm phần trăm, đủ nhỏ để một báo cáo chỉ nhìn Accuracy
kết luận sai rằng ``mô hình học được chẳng hơn bao nhiêu''. $F_1$ nói đúng sự
thật: $0{,}0000$ so với $0{,}5895$.
\end{phathien}

\section{Thực nghiệm 1 --- Đối đầu năm mô hình}
\label{sec:tn1-dia}

\textbf{Câu hỏi.} Trong năm nguyên lý học khác nhau, nguyên lý nào hợp với cấu
trúc của dữ liệu y sinh này?

\textbf{Thiết kế.} Năm mô hình nhận \emph{cùng} một biểu diễn, \emph{cùng} một
phép chia dữ liệu, \emph{cùng} một hạt giống ngẫu nhiên. Chỉ thuật toán thay đổi.

\begin{table}[H]
\centering
\caption{Kết quả năm mô hình phân lớp trên tập kiểm tra, sắp theo $F_1$}
\label{tab:ketqua-dia}
\small
\begin{tabular}{@{}lrrrrr@{}}
\toprule
\textbf{Mô hình} & \textbf{Accuracy} & \textbf{Precision} & \textbf{Recall} & \textbf{$F_1$} & \textbf{ROC-AUC} \\
\midrule
Mô hình cơ sở        & 0{,}6547 & 0{,}0000 & 0{,}0000 & 0{,}0000 & --- \\
\midrule
K-Nearest Neighbors  & 0{,}7842 & 0{,}6957 & \textbf{0{,}6667} & \textbf{0{,}6809} & 0{,}8462 \\
Logistic Regression  & \textbf{0{,}7914} & \textbf{0{,}7436} & 0{,}6042 & 0{,}6667 & \textbf{0{,}8709} \\
XGBoost              & 0{,}7194 & 0{,}5957 & 0{,}5833 & 0{,}5895 & 0{,}7905 \\
Random Forest        & 0{,}7194 & 0{,}6154 & 0{,}5000 & 0{,}5517 & 0{,}8191 \\
Decision Tree        & 0{,}7266 & 0{,}6786 & 0{,}3958 & 0{,}5000 & 0{,}7446 \\
\bottomrule
\end{tabular}
\end{table}

\begin{figure}[H]
\centering
\includegraphics[width=\textwidth]{dia_model_comparison.png}
\caption[So sánh năm mô hình phân lớp trên bốn độ đo]%
{So sánh năm mô hình phân lớp trên cùng tập kiểm tra, bốn độ đo}
\label{fig:dia-comparison}
\end{figure}

\bieudo{Biểu đồ cột nhóm (grouped bar chart). Mỗi nhóm là một mô hình, bốn cột
trong nhóm là bốn độ đo.}

\nhanxet{Ba điều đáng chú ý.

\textbf{(1) Cả năm mô hình đều vượt mô hình cơ sở} ở $F_1$ và Recall --- việc
học mang lại giá trị thật, không phải chỉ là trang trí.

\textbf{(2) Hai mô hình đơn giản nhất thắng.} K-Nearest Neighbors và Logistic
Regression đứng đầu, còn hai mô hình tổ hợp phức tạp hơn hẳn --- Random Forest
và XGBoost --- lại xếp sau. Với 552 mẫu huấn luyện và 12 đặc trưng, các mô hình
có phương sai cao không đủ dữ liệu để phát huy; đây là minh hoạ sách giáo khoa
cho đánh đổi độ chệch -- phương sai.

\textbf{(3) Không có mô hình nào thắng trên mọi độ đo.} Logistic Regression dẫn
đầu Accuracy, Precision và ROC-AUC; K-Nearest Neighbors dẫn đầu Recall và $F_1$.
Với mục tiêu sàng lọc, KNN là lựa chọn đúng --- nó bỏ sót ít người bệnh hơn.}

\subsection{Đường cong ROC --- khả năng xếp hạng độc lập với ngưỡng}

\begin{figure}[H]
\centering
\includegraphics[width=0.72\textwidth]{dia_roc_curves.png}
\caption[Đường cong ROC của năm mô hình phân lớp]%
{Đường cong ROC của năm mô hình phân lớp trên tập kiểm tra, kèm diện tích dưới
đường cong (AUC)}
\label{fig:dia-roc}
\end{figure}

\bieudo{Đường cong ROC (Receiver Operating Characteristic). Trục ngang là tỷ lệ
Dương tính giả, trục dọc là Recall; mỗi điểm trên đường cong ứng với một ngưỡng
quyết định. Đường chéo nét đứt là mốc đoán ngẫu nhiên.}

\phanphoi{Cả năm đường đều nằm hẳn phía trên đường chéo --- mọi mô hình đều xếp
hạng tốt hơn ngẫu nhiên. Chúng bó lại thành một chùm hẹp ở góc trên bên trái rồi
toả ra ở vùng giữa.}

\chitiet{Logistic Regression dẫn đầu với AUC $=0{,}8709$, bám sát là K-Nearest
Neighbors $0{,}8462$ và Random Forest $0{,}8191$. XGBoost đạt $0{,}7905$, còn
Decision Tree tụt hẳn xuống $0{,}7446$ --- đường cong của nó có dạng gãy khúc
thành vài đoạn thẳng dài thay vì cong mượt.}

\begin{phathien}{Thứ tự trên ROC không trùng thứ tự trên $F_1$ --- và đó không phải mâu thuẫn}
Logistic Regression dẫn đầu AUC nhưng chỉ đứng thứ hai về $F_1$; K-Nearest
Neighbors thì ngược lại. Hai độ đo trả lời hai câu hỏi khác nhau:

\begin{itemize}
  \item \textbf{AUC} hỏi \emph{mô hình có xếp người bệnh lên trên người khoẻ
        không}, xét trên mọi ngưỡng cùng lúc.
  \item \textbf{$F_1$} hỏi \emph{tại ngưỡng $0{,}50$, mô hình cắt có tốt không}.
\end{itemize}

Một mô hình xếp hạng giỏi mà bị cắt sai chỗ vẫn cho $F_1$ thấp. Đó chính là khe
hở mà Thực nghiệm~2 khai thác ở Mục~\ref{sec:tn2-dia}: giữ nguyên mô hình, chỉ
đổi điểm cắt, và $F_1$ tăng trong khi AUC không đổi một chút nào.

Dạng gãy khúc của đường cong Decision Tree cũng là một tín hiệu chẩn đoán: cây
sâu 6 tầng chỉ trả về vài giá trị xác suất rời rạc, nên đường cong chỉ có vài
điểm gấp. Chi tiết này sẽ quay lại ở Mục~\ref{sec:tn2-dia}.
\end{phathien}

\subsection{Lưới ma trận nhầm lẫn --- sai theo kiểu gì}

\begin{figure}[H]
\centering
\includegraphics[width=\textwidth]{dia_confusion_grid.png}
\caption[Ma trận nhầm lẫn của năm mô hình trên cùng một thang màu]%
{Ma trận nhầm lẫn của năm mô hình phân lớp, dùng chung một thang màu để so sánh
trực tiếp}
\label{fig:dia-cm-grid}
\end{figure}

\bieudo{Lưới bản đồ nhiệt (heatmap grid), mỗi ô là một ma trận nhầm lẫn $2\times2$.
Mọi ô dùng chung một thang màu nên độ đậm so sánh được trực tiếp giữa các mô
hình. Số FN và FP ghi ngay dưới tên mô hình.}

\phanphoi{Cấu trúc giống nhau ở cả năm ô: ô góc trên trái (TN) luôn đậm nhất,
ba ô còn lại nhạt hơn hẳn --- hệ quả của việc lớp âm tính đông gần gấp đôi.
Điều thay đổi giữa các mô hình là \emph{tỷ lệ} giữa hai ô sai.}

\chitiet{K-Nearest Neighbors có FN thấp nhất ($16$) --- nó bỏ sót ít người bệnh
nhất. Decision Tree có FN cao nhất ($29$), gần gấp đôi, dù Accuracy của nó
($0{,}7266$) còn cao hơn XGBoost ($0{,}7194$). Ở chiều ngược lại, Decision Tree
lại có FP thấp nhất ($9$): nó thận trọng tới mức hiếm khi báo động, và cái giá
phải trả là bỏ sót.}

\nhanxet{Đây là bằng chứng trực quan cho luận điểm ở Mục~\ref{sec:baseline}:
\textbf{Accuracy che mất đúng loại sai lầm đắt nhất}. Hai mô hình chênh nhau
$0{,}7$ điểm phần trăm Accuracy nhưng chênh nhau $13$ người bệnh bị bỏ sót trên
tổng số $48$ --- tức hơn một phần tư số ca cần tìm.}

\section{Thực nghiệm 2 --- Siêu tham số và ngưỡng quyết định}
\label{sec:tn2-dia}

\subsection{Khảo sát số láng giềng $k$}

\begin{figure}[H]
\centering
\includegraphics[width=0.88\textwidth]{dia_knn_k_sweep.png}
\caption[Ảnh hưởng của số láng giềng $k$ tới bốn độ đo]%
{Ảnh hưởng của số láng giềng $k$ tới bốn độ đo của K-Nearest Neighbors}
\label{fig:dia-k-sweep}
\end{figure}

\bieudo{Biểu đồ đường đa tuyến (multi-line plot). Trục ngang là giá trị $k$,
mỗi đường là một độ đo.}

\phanphoi{Dạng chữ U ngược đúng như lý thuyết đánh đổi độ chệch -- phương sai
dự báo, với một đỉnh rõ ràng ở đầu dải và một vùng bằng phẳng kéo dài về sau.}

\chitiet{Tại $k=1$, cả bốn độ đo đều thấp nhất ($F_1=0{,}4222$) --- mô hình học
thuộc từng điểm huấn luyện, phương sai cực đại. $F_1$ tăng vọt lên $0{,}6596$ ở
$k=3$ rồi đạt đỉnh $0{,}7292$ tại $k=5$. Từ $k=7$ trở đi đường cong đi ngang
trong dải $0{,}6667$--$0{,}6882$ mà không sụt hẳn.}

\begin{table}[H]
\centering
\caption{Khảo sát số láng giềng $k$ --- kết quả đầy đủ}
\label{tab:k-sweep}
\small
\begin{tabular}{@{}rrrrr@{}}
\toprule
\textbf{$k$} & \textbf{Accuracy} & \textbf{Precision} & \textbf{Recall} & \textbf{$F_1$} \\
\midrule
1  & 0{,}6259 & 0{,}4524 & 0{,}3958 & 0{,}4222 \\
3  & 0{,}7698 & 0{,}6739 & 0{,}6458 & 0{,}6596 \\
\textbf{5}  & \textbf{0{,}8129} & \textbf{0{,}7292} & \textbf{0{,}7292} & \textbf{0{,}7292} \\
7  & 0{,}7842 & 0{,}6875 & 0{,}6875 & 0{,}6875 \\
9  & 0{,}7770 & 0{,}6667 & 0{,}7083 & 0{,}6869 \\
11 & 0{,}7842 & 0{,}6957 & 0{,}6667 & 0{,}6809 \\
15 & 0{,}7914 & 0{,}7209 & 0{,}6458 & 0{,}6813 \\
21 & 0{,}7842 & 0{,}7143 & 0{,}6250 & 0{,}6667 \\
31 & 0{,}7914 & 0{,}7111 & 0{,}6667 & 0{,}6882 \\
\bottomrule
\end{tabular}
\end{table}

\nhanxet{Giá trị $k=11$ dùng trong danh mục mô hình ở Thực nghiệm~1 là quy tắc
kinh nghiệm $k\approx\sqrt{n}$; khảo sát cho thấy $k=5$ tốt hơn ($F_1$ $0{,}7292$
so với $0{,}6809$). Chênh lệch $0{,}048$ $F_1$ tương ứng khoảng 3 ca trên 139 ---
nằm trong biên độ nhiễu lấy mẫu của một tập kiểm tra cỡ này, nên đây là một cải
thiện \emph{có hướng} chứ chưa phải một khẳng định chắc chắn.}

\subsection{Quét ngưỡng quyết định}

Đây là thực nghiệm quan trọng nhất của cả chương, vì nó thay đổi hệ thống mà
\emph{không} đụng tới mô hình.

\begin{figure}[H]
\centering
\includegraphics[width=\textwidth]{dia_threshold_curve.png}
\caption[Ảnh hưởng của ngưỡng quyết định $\theta$ tới bốn độ đo]%
{Ảnh hưởng của ngưỡng quyết định $\theta$ tới bốn độ đo, mô hình XGBoost}
\label{fig:dia-threshold}
\end{figure}

\bieudo{Biểu đồ đường đa tuyến. Trục ngang là ngưỡng $\theta$ quét từ $0{,}60$
xuống $0{,}05$; đường nét đứt dọc đánh dấu ngưỡng tối ưu theo $F_1$.}

\phanphoi{Hai đường chính chạy ngược chiều nhau và cắt nhau --- dấu hiệu kinh
điển của một đánh đổi. Hai đường còn lại biến thiên nhẹ.}

\chitiet{Recall tăng đơn điệu khi $\theta$ giảm: từ $0{,}4583$ tại
$\theta=0{,}60$ lên $0{,}9167$ tại $\theta=0{,}05$. Precision đi ngược lại, từ
$0{,}5946$ xuống $0{,}4583$. Hai đường cắt nhau quanh $\theta=0{,}50$. $F_1$ ---
trung bình điều hoà của hai đường ấy --- đạt đỉnh $0{,}6508$ tại
$\theta=0{,}15$. Accuracy gần như đi ngang trong dải $0{,}60$--$0{,}72$ suốt cả
quá trình, một lần nữa cho thấy nó không nhạy với điều đang thay đổi.}

\begin{table}[H]
\centering
\caption{Quét ngưỡng quyết định --- bốn ô ma trận nhầm lẫn tại mỗi ngưỡng}
\label{tab:threshold-dia}
\small
\begin{tabular}{@{}rrrrrrrrr@{}}
\toprule
\textbf{$\theta$} & \textbf{Accuracy} & \textbf{Precision} & \textbf{Recall} & \textbf{$F_1$} & \textbf{TP} & \textbf{FN} & \textbf{FP} & \textbf{TN} \\
\midrule
0{,}60 & 0{,}7050 & 0{,}5946 & 0{,}4583 & 0{,}5176 & 22 & 26 & 15 & 76 \\
0{,}50 & 0{,}7194 & 0{,}5957 & 0{,}5833 & 0{,}5895 & 28 & 20 & 19 & 72 \\
0{,}40 & 0{,}7194 & 0{,}5789 & 0{,}6875 & 0{,}6286 & 33 & 15 & 24 & 67 \\
0{,}30 & 0{,}7050 & 0{,}5538 & 0{,}7500 & 0{,}6372 & 36 & 12 & 29 & 62 \\
0{,}25 & 0{,}7122 & 0{,}5606 & 0{,}7708 & 0{,}6491 & 37 & 11 & 29 & 62 \\
0{,}20 & 0{,}6978 & 0{,}5417 & 0{,}8125 & 0{,}6500 & 39 &  9 & 33 & 58 \\
\textbf{0{,}15} & 0{,}6835 & 0{,}5256 & \textbf{0{,}8542} & \textbf{0{,}6508} & \textbf{41} & \textbf{7} & 37 & 54 \\
0{,}10 & 0{,}6403 & 0{,}4886 & 0{,}8958 & 0{,}6324 & 43 &  5 & 45 & 46 \\
0{,}05 & 0{,}5971 & 0{,}4583 & 0{,}9167 & 0{,}6111 & 44 &  4 & 52 & 39 \\
\bottomrule
\end{tabular}
\end{table}

\begin{figure}[H]
\centering
\begin{subfigure}[b]{0.485\textwidth}
  \includegraphics[width=\textwidth]{dia_confusion_default.png}
  \caption{Ngưỡng mặc định $\theta=0{,}50$}
  \label{fig:dia-cm-default}
\end{subfigure}
\hfill
\begin{subfigure}[b]{0.485\textwidth}
  \includegraphics[width=\textwidth]{dia_confusion_tuned.png}
  \caption{Ngưỡng tối ưu $\theta=0{,}15$}
  \label{fig:dia-cm-tuned}
\end{subfigure}
\caption[Ma trận nhầm lẫn trước và sau khi hạ ngưỡng quyết định]%
{Ma trận nhầm lẫn của cùng một mô hình XGBoost, trước và sau khi hạ ngưỡng
quyết định}
\label{fig:dia-cm-pair}
\end{figure}

\nhanxet{Đọc thẳng từ Bảng~\ref{tab:threshold-dia}: hạ ngưỡng từ $0{,}50$ xuống
$0{,}15$ đưa số người bệnh bị bỏ sót từ \textbf{20 xuống 7} --- cứu được 13 ca
trên tổng số 48. Cái giá là số báo động nhầm tăng từ 19 lên 37, tức thêm 18
người khoẻ phải đi xét nghiệm lại.

Trong bài toán sàng lọc, đây là một đổi chác rõ ràng có lợi: một lần xét nghiệm
thừa tốn vài trăm nghìn đồng, còn một ca đái tháo đường bị bỏ sót có thể dẫn tới
biến chứng thận, mắt hoặc tim mạch sau nhiều năm. \textbf{Điểm quan trọng nhất:
mô hình không hề thay đổi} --- cùng bộ trọng số, cùng ROC-AUC $0{,}7905$. Thứ
thay đổi là chỗ ta cắt.}

\begin{phathien}{Hai bẫy phương pháp gặp phải trong chính thực nghiệm này}
\textbf{Bẫy 1 --- quét ngưỡng cần mô hình cho xác suất liên tục.} Bản chạy đầu
tiên thực hiện trên Decision Tree sâu 6 tầng. Bảy ngưỡng từ $0{,}50$ xuống
$0{,}20$ cho ra \emph{đúng một} kết quả giống hệt nhau. Nguyên nhân: cây nông
chỉ trả về vài giá trị xác suất rời rạc, nên phần lớn ngưỡng rơi vào cùng một
khoảng và không đổi được nhãn nào --- đúng hiện tượng đã thấy qua dạng gãy khúc
của đường ROC ở Hình~\ref{fig:dia-roc}. Chuyển sang XGBoost mới thấy được hiệu
ứng thật, với 12 mức kết quả khác nhau trên 13 ngưỡng.

\textbf{Bẫy 2 --- cực trị chạm biên dải khảo sát thì chưa phải cực trị.} Dải
quét ban đầu dừng ở $\theta=0{,}20$ và $F_1$ đạt đỉnh đúng tại đầu mút ấy. Chỉ
sau khi nới dải xuống $0{,}05$ mới xác nhận được đỉnh thật nằm ở $0{,}15$. Bẫy
này lặp lại lần thứ hai ở bài toán bất động sản với độ sâu rừng ngẫu nhiên ---
xem Mục~\ref{sec:tn2-hou}.
\end{phathien}

\section{Thực nghiệm 3 --- Ảnh hưởng của cách biểu diễn dữ liệu}
\label{sec:tn3-dia}

\textbf{Câu hỏi.} Giữ nguyên thuật toán, chỉ đổi cách biểu diễn --- kết quả thay
đổi bao nhiêu, và so với biên độ do đổi thuật toán thì thế nào?

\textbf{Thiết kế.} Bốn phương án biểu diễn, cùng chạy trên Logistic Regression,
cùng phép chia dữ liệu. Bốn phương án phản ánh đúng ba tầng chuyển hoá đã định
nghĩa ở Mục~\ref{sec:ba-tang-bieudien}.

\begin{table}[H]
\centering
\caption{Bốn phương án biểu diễn dữ liệu và kết quả trên Logistic Regression}
\label{tab:rep-dia}
\small
\begin{tabular}{@{}lrrrr@{}}
\toprule
\textbf{Phương án biểu diễn} & \textbf{Accuracy} & \textbf{Precision} & \textbf{Recall} & \textbf{$F_1$} \\
\midrule
Làm sạch, chưa chuẩn hoá                    & \textbf{0{,}8058} & 0{,}7561 & \textbf{0{,}6458} & \textbf{0{,}6966} \\
Thô --- giữ nguyên giá trị 0 phi lý         & 0{,}7986 & 0{,}7500 & 0{,}6250 & 0{,}6818 \\
Làm sạch $+$ chuẩn hoá                      & 0{,}7986 & 0{,}7500 & 0{,}6250 & 0{,}6818 \\
Làm sạch $+$ chuẩn hoá $+$ đặc trưng thiết kế & 0{,}7914 & 0{,}7436 & 0{,}6042 & 0{,}6667 \\
\bottomrule
\end{tabular}
\end{table}

\begin{figure}[H]
\centering
\includegraphics[width=\textwidth]{dia_representation.png}
\caption[Ảnh hưởng của cách biểu diễn dữ liệu tới $F_1$]%
{Ảnh hưởng của cách biểu diễn dữ liệu tới $F_1$ của Logistic Regression}
\label{fig:dia-representation}
\end{figure}

\bieudo{Biểu đồ cột ngang, sắp giảm dần theo $F_1$, giá trị ghi ở đầu mỗi cột.}

\phanphoi{Bốn cột gần bằng nhau. Khoảng cách giữa cột cao nhất và cột thấp nhất
nhỏ tới mức phải đọc con số mới thấy.}

\chitiet{$F_1$ trải từ $0{,}6667$ tới $0{,}6966$ --- biên độ $0{,}0300$. Hai
phương án giữa cho kết quả \emph{trùng khít nhau tới bốn chữ số thập phân}
($0{,}6818$), tức việc chuẩn hoá không đổi được một ca nào.}

\begin{phathien}{Kết quả đi ngược giả thuyết ban đầu --- và đây là phát hiện chính}
\label{sec:dactrung-tot}
Giả thuyết trước khi chạy là biểu diễn tốt hơn sẽ cho kết quả tốt hơn, với
phương án đầy đủ nhất đứng đầu. Kết quả \textbf{ngược lại}: phương án đầy đủ
nhất --- có cả chuẩn hoá lẫn bốn đặc trưng thiết kế --- xếp \textbf{cuối cùng}.

Ba lời giải thích, mỗi lời giải thích kiểm chứng được:

\begin{enumerate}
  \item \textbf{Chuẩn hoá không giúp gì cho Logistic Regression ở đây.} Hai
        phương án ``có'' và ``không có'' chuẩn hoá cho kết quả trùng khít. Điều
        này đúng lý thuyết: hồi quy logistic không dựa trên khoảng cách, phép
        chuẩn hoá tuyến tính chỉ đổi thang trọng số chứ không đổi ranh giới
        quyết định. Nó \emph{bắt buộc} với K-Nearest Neighbors, nhưng vô hại và
        vô ích với mô hình này.
  \item \textbf{Đặc trưng thiết kế \emph{làm hại} Logistic Regression.} $F_1$
        giảm từ $0{,}6818$ xuống $0{,}6667$. Nguyên nhân là đa cộng tuyến dưới
        phạt $L_2$: \texttt{Glucose\_BMI\_Risk} tương quan $0{,}82$ với
        \texttt{Glucose} và $0{,}73$ với \texttt{BMI}
        (Hình~\ref{fig:dia-corr}), nên số hạng
        $\lambda\lVert\bm\theta\rVert_2^2$ ở~\eqref{eq:regularized} phân tán
        trọng số ra ba cột gần trùng nhau thay vì dồn vào một cột.
  \item \textbf{Toàn bộ biên độ nằm trong nhiễu lấy mẫu.} $0{,}0300$ $F_1$ trên
        139 ca kiểm tra tương ứng chênh lệch chỉ \textbf{1--9 ca} giữa các
        phương án. Không phương án nào tách khỏi phần còn lại một cách có ý
        nghĩa thống kê.
\end{enumerate}

\textbf{Hệ quả khái quát:} ``đặc trưng tốt'' không phải thuộc tính của riêng đặc
trưng, mà là thuộc tính của \emph{cặp đặc trưng -- thuật toán}. Chính
\texttt{Glucose\_BMI\_Risk} vừa là đặc trưng tương quan mạnh nhất với nhãn
($0{,}5199$) vừa là đặc trưng quan trọng nhất của XGBoost ($0{,}1874$, xem
Hình~\ref{fig:dia-importance}) --- nhưng lại làm Logistic Regression kém đi.
\end{phathien}

\subsection{So sánh biên độ: biểu diễn hay thuật toán?}

\begin{table}[H]
\centering
\caption{Biên độ do đổi biểu diễn so với biên độ do đổi thuật toán, bài toán tiểu đường}
\label{tab:biendo-dia}
\small
\begin{tabular}{@{}lrl@{}}
\toprule
\textbf{Nguồn biến thiên} & \textbf{Biên độ $F_1$} & \textbf{Ghi chú} \\
\midrule
Đổi cách biểu diễn (4 phương án) & 0{,}0300 & $0{,}6667 \to 0{,}6966$ \\
Đổi thuật toán (5 mô hình)       & \textbf{0{,}1809} & $0{,}5000 \to 0{,}6809$ \\
\midrule
\textbf{Tỷ số} & \multicolumn{2}{l}{\textbf{Thuật toán chi phối, gấp 6 lần}} \\
\bottomrule
\end{tabular}
\end{table}

Kết luận này sẽ được đối chiếu với kết quả \emph{trái ngược} ở bài toán bất động
sản tại Mục~\ref{sec:phanchieu}.

\section{Thực nghiệm 4 --- Kiểm chuẩn ngoài phân phối}
\label{sec:tn4-dia}

\textbf{Câu hỏi.} Điểm số trên tập kiểm tra có lạc quan quá không?

\begin{figure}[H]
\centering
\includegraphics[width=0.92\textwidth]{dia_ood_comparison.png}
\caption[So sánh $F_1$ giữa tập kiểm tra và tập ngoại kiểm]%
{So sánh $F_1$ của năm mô hình giữa tập kiểm tra và tập ngoại kiểm}
\label{fig:dia-ood}
\end{figure}

\begin{table}[H]
\centering
\caption{Kết quả trên tập ngoại kiểm (77 ca chưa từng dùng trong phát triển)}
\label{tab:ood-dia}
\small
\begin{tabular}{@{}lrrrrr@{}}
\toprule
\textbf{Mô hình} & \textbf{Accuracy} & \textbf{Precision} & \textbf{Recall} & \textbf{$F_1$} & \textbf{ROC-AUC} \\
\midrule
XGBoost             & \textbf{0{,}7922} & 0{,}6897 & \textbf{0{,}7407} & \textbf{0{,}7143} & 0{,}8800 \\
Random Forest       & \textbf{0{,}7922} & 0{,}7037 & 0{,}7037 & 0{,}7037 & \textbf{0{,}8904} \\
K-Nearest Neighbors & \textbf{0{,}7922} & \textbf{0{,}7391} & 0{,}6296 & 0{,}6800 & \textbf{0{,}8904} \\
Logistic Regression & 0{,}7403 & 0{,}6400 & 0{,}5926 & 0{,}6154 & 0{,}8511 \\
Decision Tree       & 0{,}7403 & 0{,}6842 & 0{,}4815 & 0{,}5652 & 0{,}7774 \\
\bottomrule
\end{tabular}
\end{table}

\nhanxet{Kết quả trên tập ngoại kiểm \textbf{cao hơn} tập kiểm tra ở hầu hết mô
hình --- XGBoost tăng $F_1$ từ $0{,}5895$ lên $0{,}7143$. Đây không phải bằng
chứng mô hình tổng quát hoá xuất sắc; nó chủ yếu là dao động lấy mẫu trên một
tập chỉ 77 ca. Điều đáng tin hơn là \emph{thứ hạng đảo lộn}: XGBoost từ hạng ba
lên hạng nhất, Logistic Regression từ hạng hai xuống hạng tư.

Hệ quả phương pháp: chênh lệch $F_1$ dưới khoảng $0{,}05$ giữa hai mô hình trên
tập cỡ này \textbf{không đủ để kết luận mô hình nào tốt hơn}. Đây cũng là lý do
Mục~\ref{sec:tn3-dia} không tuyên bố phương án biểu diễn nào thắng.}

\section{Mức độ quan trọng của đặc trưng}
\label{sec:importance-dia}

\begin{figure}[H]
\centering
\includegraphics[width=0.88\textwidth]{dia_feature_importance.png}
\caption[Mức độ quan trọng của đặc trưng theo XGBoost]%
{Mức độ quan trọng của đặc trưng theo mô hình XGBoost}
\label{fig:dia-importance}
\end{figure}

\bieudo{Biểu đồ cột ngang, sắp tăng dần từ trên xuống. Độ quan trọng đo bằng
mức giảm độ vẩn tích luỹ mà mỗi đặc trưng mang lại qua toàn bộ cây.}

\phanphoi{Phân bố khá đều, không có đặc trưng nào áp đảo. Hai cột đầu tách khỏi
phần còn lại, chín cột sau tạo thành một dải liên tục.}

\chitiet{\texttt{Glucose\_BMI\_Risk} dẫn đầu với $0{,}1874$, kế đến
\texttt{Glucose} gốc $0{,}1233$. Nhóm giữa gồm \texttt{Age} $0{,}0940$,
\texttt{Age\_Group} $0{,}0815$, \texttt{BMI} $0{,}0778$ và \texttt{Pregnancies}
$0{,}0776$. Cuối bảng là \texttt{Glucose\_Level} với $0{,}0330$.}

\nhanxet{Ba quan sát.

\textbf{(1) Đặc trưng thiết kế đứng đầu} --- bằng chứng thứ hai (sau
Bảng~\ref{tab:corr-dia}) rằng tri thức y học đưa vào đúng chỗ có giá trị đo
được, ít nhất là với mô hình cây.

\textbf{(2) Rời rạc hoá làm mất thông tin với mô hình cây.}
\texttt{Glucose\_Level} xếp cuối trong khi \texttt{Glucose} gốc xếp thứ hai.
Điều này hợp lý: cây quyết định \emph{tự} tìm ngưỡng chia tối ưu trên biến liên
tục, nên việc rời rạc hoá thủ công thành ba mức chỉ giới hạn không gian tìm kiếm
của nó. Rời rạc hoá có ích cho mô hình tuyến tính, không có ích cho cây.

\textbf{(3) Không có đặc trưng nào vượt $0{,}19$} --- tín hiệu phân tán mỏng
trên nhiều biến. Điều này giải thích vì sao trần hiệu năng trên tập Pima nằm ở
mức $0{,}77$--$0{,}79$ Accuracy trong toàn bộ tài liệu công bố, và vì sao mọi
kết quả vượt xa mốc ấy đều đáng nghi (xem lại khung ở Mục~\ref{sec:roro}).}

\section{Tích hợp Đồ thị Tri thức Y sinh}
\label{sec:kg-dia}

\subsection{Từ xác suất tới tầng nguy cơ}

Ngưỡng phân tầng lấy trực tiếp từ Thực nghiệm~2, không phải chọn tuỳ ý:

\begin{table}[H]
\centering
\caption{Ba tầng nguy cơ lâm sàng và ánh xạ sang mã bệnh quốc tế}
\label{tab:tang-dia}
\small
\begin{tabularx}{\textwidth}{@{}L{3.6cm}L{1.7cm}L{1.7cm}X@{}}
\toprule
\textbf{Tầng nguy cơ} & \textbf{Ngưỡng $\hat p$} & \textbf{ICD-10} & \textbf{Tên mã bệnh} \\
\midrule
Nhóm Nguy cơ Cao        & $\ge 0{,}25$ & E11   & Đái tháo đường Type 2 \\
Nhóm Tiền Đái tháo đường & $\ge 0{,}15$ & R73.0 & Tăng glucose máu bất thường \\
Nhóm Khoẻ mạnh          & $\ge 0{,}00$ & Z13.1 & Khám sàng lọc đái tháo đường \\
\bottomrule
\end{tabularx}
\end{table}

\subsection{Đồ thị của một ca cụ thể}

\begin{figure}[H]
\centering
\includegraphics[width=\textwidth]{kg_medical.png}
\caption[Đồ thị Tri thức Y sinh của bệnh nhân \#192]%
{Đồ thị Tri thức Y sinh của bệnh nhân \#192 --- Nhóm Nguy cơ Cao, 22 thực thể
và 31 cạnh}
\label{fig:kg-medical}
\end{figure}

\bieudo{Đồ thị có hướng vẽ theo bố cục phân tầng (layered graph layout). Màu nút
mã hoá loại thực thể theo chú giải bên phải, nhãn trên cạnh là tên quan hệ.}

\phanphoi{Cấu trúc hình quạt mở dần từ trái sang phải: một nút gốc duy nhất
(Bệnh nhân) phân nhánh qua các tầng suy luận rồi mở rộng ra tầng hàng hoá và
dịch vụ.}

\chitiet{Tầng 1 bên trái gồm bốn chỉ số sinh học cùng đặc trưng tương tác
$\text{Glucose}\times\text{BMI}=4.834$. Tầng 2 ở giữa là chuỗi Mô hình
K-Nearest Neighbors $\to$ Xác suất $54{,}5\%$ $\to$ ICD-10 E11 $\to$ Nhóm Nguy
cơ Cao. Tầng 3 toả ra ba nhóm: hai thiết bị, ba mặt hàng dinh dưỡng, ba hướng
dẫn lâm sàng. Tầng 4 gom lại thành một nút Gói Theo dõi Chuyên sâu. Tầng 5 bên
phải hội tụ về nút Nhà thuốc FPT Long Châu cùng hai dịch vụ tại quầy và một
đường chuyển tuyến chuyên khoa.}

\subsection{Chuỗi suy luận đầy đủ}

\begin{table}[H]
\centering
\caption{Chuỗi bộ ba tri thức của bệnh nhân \#192, từ số đo tới nơi mua hàng}
\label{tab:triplet-dia}
\scriptsize
\begin{tabularx}{\textwidth}{@{}L{4.1cm}L{3.5cm}X@{}}
\toprule
\textbf{Chủ thể} & \textbf{Quan hệ} & \textbf{Đối tượng} \\
\midrule
Bệnh nhân & \texttt{HAS\_BIOMARKER} & Glucose $=159$ mg/dL \\
Bệnh nhân & \texttt{HAS\_BIOMARKER} & BMI $=30{,}4$ kg/m$^2$ \\
Bệnh nhân & \texttt{CALCULATED\_RISK} & Glucose $\times$ BMI $=4.834$ \\
Bệnh nhân & \texttt{EVALUATED\_BY} & Mô hình K-Nearest Neighbors \\
Mô hình K-Nearest Neighbors & \texttt{PRODUCES} & Xác suất rủi ro $54{,}5\%$ \\
Xác suất rủi ro $54{,}5\%$ & \texttt{MAPS\_TO\_DISEASE} & ICD-10 E11 \\
Bệnh nhân & \texttt{STRATIFIED\_AS} & Nhóm Nguy cơ Cao \\
ICD-10 E11 & \texttt{REQUIRES\_DEVICE} & Máy đo đường huyết Accu-Chek Instant \\
ICD-10 E11 & \texttt{RECOMMENDS\_NUTRITION} & Sữa Abbott Glucerna 850g \\
Nhóm Nguy cơ Cao & \texttt{ADVISES} & Khám chuyên khoa Nội tiết \\
Nhóm Nguy cơ Cao & \texttt{RECOMMENDS\_PACKAGE} & \textbf{Gói Theo dõi Chuyên sâu} \\
Gói Theo dõi Chuyên sâu & \texttt{INCLUDES\_PRODUCT} & (5 mặt hàng) \\
Gói Theo dõi Chuyên sâu & \texttt{FULFILLED\_BY} & \textbf{Nhà thuốc FPT Long Châu} \\
Nhà thuốc FPT Long Châu & \texttt{PROVIDES\_SERVICE} & Đo đường huyết miễn phí tại quầy \\
ICD-10 E11 & \texttt{REFERS\_TO} & Khám Bác sĩ chuyên khoa Nội tiết \\
\bottomrule
\end{tabularx}
\end{table}

\subsection{Gói chăm sóc sức khoẻ và giải pháp dược phẩm}
\label{sec:goi-cham-soc}

Đây là điểm mà hệ thống trả lời được câu hỏi cuối cùng của người dùng:
\emph{vậy tôi phải mua gì, và hết bao nhiêu tiền}.

\begin{table}[H]
\centering
\caption{Gói Theo dõi Chuyên sâu --- danh mục mặt hàng cho Nhóm Nguy cơ Cao}
\label{tab:goi-cao}
\scriptsize
\begin{tabularx}{\textwidth}{@{}L{1.9cm}XL{2.6cm}L{1.35cm}@{}}
\toprule
\textbf{Mã SKU} & \textbf{Mặt hàng} & \textbf{Khoảng giá} & \textbf{Định kỳ} \\
\midrule
LC-DEV-001 & Máy đo đường huyết Accu-Chek Instant        & 790.000 -- 1.090.000 đ & Không \\
LC-DEV-002 & Que thử đường huyết Accu-Chek (hộp 50 que)  & 300.000 -- 420.000 đ   & Có \\
LC-DEV-003 & Kim chích máu Accu-Chek Softclix (hộp 25)   & 90.000 -- 150.000 đ    & Có \\
LC-NUT-001 & Sữa Abbott Glucerna 850g                    & 690.000 -- 890.000 đ   & Có \\
LC-NUT-003 & Viên uống Dây thìa canh Diabetna (hộp 40)   & 120.000 -- 195.000 đ   & Có \\
\midrule
\multicolumn{2}{@{}l}{\textbf{Chi phí ban đầu}}          & \textbf{1.990.000 -- 2.745.000 đ} & \\
\multicolumn{2}{@{}l}{\textbf{Duy trì hằng tháng}}       & \textbf{1.200.000 -- 1.655.000 đ} & \\
\multicolumn{2}{@{}l}{\textbf{Tổng 12 tháng đầu}}        & \textbf{15.190.000 -- 20.950.000 đ} & \\
\bottomrule
\end{tabularx}
\end{table}

\begin{table}[H]
\centering
\caption{Ba gói chăm sóc theo tầng nguy cơ}
\label{tab:ba-goi}
\small
\begin{tabularx}{\textwidth}{@{}L{3.9cm}rL{2.8cm}X@{}}
\toprule
\textbf{Gói} & \textbf{Số mặt hàng} & \textbf{Chi phí ban đầu} & \textbf{Tần suất theo dõi} \\
\midrule
Theo dõi Chuyên sâu       & 5 & 1{,}99 -- 2{,}75 tr.đ & Đo đường huyết đói và sau ăn 2 giờ, mỗi ngày \\
Dự phòng Tiền Đái tháo đường & 4 & 1{,}43 -- 2{,}08 tr.đ & Kiểm tra lại đường huyết sau 3 tháng \\
Duy trì Sức khoẻ          & 1 & 0{,}15 -- 0{,}32 tr.đ & Khám sức khoẻ định kỳ hằng năm \\
\bottomrule
\end{tabularx}
\end{table}

\nhanxet{Chi phí ban đầu và chi phí duy trì được tách riêng có chủ đích. Gói
Theo dõi Chuyên sâu có máy đo là khoản mua một lần trị giá tới $1{,}09$ triệu
đồng; gộp chung hai con số sẽ khiến gói trông đắt gấp rưỡi thực tế của tháng thứ
hai trở đi. Đây là loại minh bạch mà một bảng khuyến nghị chỉ liệt kê tên hàng
không thể cung cấp.

Về nguồn số liệu, xin nhắc lại ranh giới đã nêu ở Mục~\ref{sec:ranh-gioi-so-lieu}:
tên mặt hàng là tên thật, giá là khoảng tham khảo chốt ngày 26/08/2026, và mỗi
mặt hàng mang một đường dẫn tra cứu giá hiện hành. Thuốc kê đơn --- Metformin và
nhóm hạ đường huyết khác --- \textbf{không} nằm trong danh mục này; đồ thị chỉ
ghi lời khuyến cáo ``tuyệt đối không tự dùng khi chưa có chỉ định bác sĩ''.}

\section{Ứng dụng web}
\label{sec:app-dia}

\begin{figure}[H]
\centering
\includegraphics[width=0.86\textwidth]{app-desktop-2-chan-doan.png}
\caption[Giao diện trang Chẩn đoán Tiểu đường trên máy tính]%
{Giao diện trang Chẩn đoán Tiểu đường trên máy tính --- biểu mẫu nhập liệu, kết
quả dự báo và đồ thị tri thức dựng tại chỗ}
\label{fig:app-dia-desktop}
\end{figure}

\bieudo{Ảnh chụp màn hình bản đã triển khai công khai, khổ máy tính $1440$ px.}

\chitiet{Bố cục hai cột. Cột trái là biểu mẫu chia ba khối --- thông tin cá
nhân, kết quả xét nghiệm, tiền sử gia đình --- với ô để trống được xử lý như dữ
liệu khuyết chứ không phải giá trị $0$. Cột phải hiển thị vòng tròn xác suất,
nhãn tầng nguy cơ, mã ICD-10 và cảnh báo chuyển tuyến. Phía dưới là bốn tab:
Khuyến nghị, Gói chăm sóc, So sánh, Vì sao.}

\nhanxet{Ứng dụng nạp lại đúng các tệp \texttt{models/*.joblib} do
\texttt{run\_pipeline.py} sinh ra, nên biểu diễn lúc dự đoán trùng khít biểu
diễn lúc huấn luyện --- toàn bộ bước điền khuyết, thiết kế đặc trưng và chuẩn
hoá đều nằm bên trong \texttt{Pipeline} được lưu.

Chi tiết đáng chú ý về mặt thiết kế: khi người dùng để trống một ô, ứng dụng
hiển thị dòng cảnh báo cho biết có bao nhiêu chỉ số đang được điền bằng trung vị
của tập huấn luyện. Điều này giữ cho người dùng biết mức độ tin cậy của con số
họ vừa nhận --- một dự báo dựa trên 8 chỉ số đo đủ khác hẳn một dự báo mà 5
trong 8 chỉ số là giá trị điền.}

\begin{figure}[H]
\centering
\includegraphics[width=0.86\textwidth]{app-desktop-4-goi-cham-soc.png}
\caption[Tab Gói chăm sóc --- giải pháp dược phẩm khuyến nghị]%
{Tab \emph{Gói chăm sóc} --- năm mặt hàng của Gói Theo dõi Chuyên sâu, mỗi thẻ
mang mã SKU, khoảng giá và đường dẫn tra giá hiện hành}
\label{fig:app-dia-goi}
\end{figure}

\bieudo{Ảnh chụp màn hình tab thứ hai của trang Chẩn đoán.}

\chitiet{Phần trên là biểu mẫu đã nạp sẵn hồ sơ bệnh nhân \#192 --- Glucose
$159$ mg/dL, BMI $30{,}4$ --- cùng kết quả $55\%$ và mã ICD-10 E11. Trong tab,
ba thẻ số cho chi phí ban đầu $1{,}99$--$2{,}75$ triệu đồng, chi phí duy trì
$1{,}20$--$1{,}66$ triệu đồng mỗi tháng, và tần suất theo dõi. Bên dưới là lưới
năm thẻ mặt hàng, mỗi thẻ ghi mã SKU, nhóm hàng, khoảng giá theo đơn vị bán,
công dụng, nhãn ``mua một lần'' hoặc ``mua định kỳ'', và một liên kết tra giá.
Cuộn tiếp xuống dưới khung hình còn có bảng tóm tắt nhà thuốc --- với con số chi
phí 12 tháng đầu $15{,}19$--$20{,}95$ triệu đồng --- và khung cảnh báo nguồn giá.}

\nhanxet{Khung cảnh báo ở cuối tab là phần không được lược bỏ: nó nêu rõ ngày
chốt giá, nói rằng đây là gợi ý mua sắm chứ \textbf{không phải đơn thuốc}, và
nhắc rằng thuốc kê đơn không nằm trong danh mục. Một giao diện bán hàng bỏ dòng
này đi sẽ biến một công cụ sàng lọc thành một quảng cáo dược phẩm.}

\begin{figure}[H]
\centering
\includegraphics[width=0.34\textwidth]{app-mobile-2-chan-doan.png}
\caption[Giao diện trang Chẩn đoán trên điện thoại]%
{Giao diện trang Chẩn đoán trên điện thoại, khổ $390$ px --- biểu mẫu xếp dọc,
thanh điều hướng tự xuống dòng}
\label{fig:app-dia-mobile}
\end{figure}
